\chapter{Онолын судалгаа}
\section{Робоконий дүрэм}
Робоконий тоглох зааланд 2 баг өрсөлдөх бөгөөд, баг тус бүрд заалын тэн хагас нь оноогдоно. Мөн оноогдсон хэсэг тус бүр нь дотроо 3 zone-д хуваагдах бөгөөд, тэдгээр нь:
\begin{itemize}
    \item Martial club --- Тулааны клуб
    \item Meihua forest --- Мэйхуа ой
    \item Arena --- Тулааны талбар
\end{itemize}
нэртэй байна.

\begin{figure}[h]
    \includegraphics[width=1\textwidth]{robocon_field.png}
    \caption{\Robocon{} тэмцээний тоглох талбай}
    \label{fig:robocon_field}
\end{figure}

\newpage

Нэг багаас 2 робот өрсөлдөх бөгөөд, үүнд:
\begin{itemize}
    \item R1 робот --- Удирдлагатай эсвэл автомат байж болно.
    \item R2 робот --- Бүрэн автомат байх бөгөөд ямар ч комманд болон удирдлага авахгүй байна.
\end{itemize}

Мөн Мэйхуа ой дээр 8 ширхэг Күнгфү судар (Kung fu scroll) байрлуулна. Үүнд:
\begin{itemize}
    \item 3 ширхэг R1-ийн судар нь R1 робот авахад зориулсан байх бөгөөд, бүх талдаа тэнцээний логотой байна.
    \item 4 ширхэг R2-ийн судар нь R2 робот авахад зориулсан байх бөгөөд, бүх талдаа жинхэнэ мэргэний ясан бичиг агуулсан байна.
    \item Үлдсэн нэг нь хуурамч судар байх ба, аль ч робот авах ёсгүй байхаар, бүх талдаа хуурамч мэргэний ясан бичигтэй байна.
\end{itemize}
Дүрэм дээр жинхэнэ болон хуурамч мэргэний ясан бичиг тус бүр 15 ширхэг байх ба тэмцээний үеэд эдгээрээс 7 ширхэг жинхэнэ, 1 ширхэг хуурмаг бичиг дурын сонгогдож, Күнгфү сударт бичигдэнэ.

\begin{figure}[h]
    \includegraphics[width=\textwidth]{robocon_kfs.png}
    \caption{\Robocon{} тэмцээний Күнгфү сударт тавигдах дүрсүүд.}
    \label{fig:robocon_kfs}
\end{figure}

\subsection{Тулааны клуб zone}
Тулааны клуб zone-оос роботууд хөдөлж эхлэх бөгөөд, роботоо бэлдэх 1 минут өгнө.

\newpage
\section{Күнгфү судар таних}

Күнгфү судар танихад зураг таних машин сургалт хэрэгтэй болно. Сургахад бодит дата одоогоор байхгүй байгаа тул симуляцийн орчинд күнгфү сударуудыг үүсгэж, датасэт үүсгэх хэрэгтэй болж байна.

Мөн күнгфү судар нь 15 үнэн болон 15 хуурмаг бичиг байгаа бөгөөд, \Robocon{} тэмцээний логотой судар гэх 3 төрлийн датасет сургах хэрэгтэй болно. Учир нь R2 робот буруу судараа ерөнхийд нь ойролцоо гэж таних магадлалтай тул буруу сударуудыг нь заавал таниулж, авахгүй байх зориулалттай байна.

\subsection{Датасет боловсруулах}

Датасет боловсруулахад Unity гэх видео тоглоом хийхэд зориулсан систем ашиглаж, симуляцийн орчин үүсгэж, датасетээ боловсруулна.

Unity нь 2 хэмжээст болон 3 хэмжээст видео тоглоом хийхэд зориулагдсан game engine юм. Unity нь өндөр чанартай бодит амьдралтай ойролцоо гэрэлтүүлэг, сүүдэр, камерийн эффект тохируулснаар, янз бүрийн гэрэлтүүлэг, сүүдэр, бүрэлзэх, г.м. тохиолдлуудыг симуляц хийж, өндөр чанартай сургалтын датасет боловсруулах чадвартай байна.

\begin{figure}[h]
    \centering
    
\includegraphics[width=0.5\textwidth]{unity.png}
    \caption{Unity Game Engine лого}
    \label{fig:unity}
\end{figure}

\subsection{Модель үүсгэх}

Симуляцийн орчин бүрдүүлэхэд тоглох талбай болон судруудын модель үүсгэх хэрэгтэй.

Тоглох талбайн хэмжээ нь дүрэм дээр өгөгдсөн байгаа бөгөөд, талбайг загварчлахдаа FreeCAD ашиглсан загварчилсан болно.

\begin{figure}[h]
    \centering
    \includegraphics[width=\textwidth]{robocon-field-blueprint.png}
    \caption{Робокон тэмцээний дүрэм дээр заасан тоглох талбайн хэмжээ}
    \label{fig:field-blueprint}
\end{figure}
\begin{figure}[h]
    \centering
    
\includegraphics[width=0.7\textwidth]{logo-freecad.png}
    \caption{FreeCAD лого}
    \label{fig:freecad}
\end{figure}

FreeCAD нь математик тооцооллоор геометр дүрсүүд ашиглан биет бүтээдэг билээ. Гэвч Unity нь зөвхөн гурвалжин дүрсээр бүрдэх торон бүтэцтэй биетийг зурах чадвартай байдаг байна. Үүнд зориулсан гурвалжин торлог бүтэц агуулах файлын формат төрөл бүр байдаг бөгөөд, өргөн хэрэглэгддэг .stl болон .fbx форматуудыг ашигласан болно.

.stl формат нь FreeCAD программын үүсгэж чадах формат бөгөөд, торон биетийн гурвалжин бүрийн өнцгүүдийн координатыг хадгалдаг формат юм.

.fbx формат нь .stl форматтай адил торон биетийн бүтцийг хадгалахаас гадна тухайн биет ямар өнгөтэй, гэрлэнд ойх чанар, барзгар болон тэгш гадаргуу г.м. материаллаг шинж чанар хадгалах чадвартай байна. Энэ форматыг FreeCAD гаргах чадваргүй тул Blender ашиглан, .stl форматтай биетэд материаллаг шинж чанар олгож, .fbx форматруу хөрвүүлсэн болно.

\begin{figure}[h]
    \centering
    \includegraphics[width=\textwidth]{field-model-freecad.png}
    \caption{FreeCAD дээр үүсгэсэн тоглох талбайн модель}
\end{figure}
\begin{figure}[h]
    \centering
    \includegraphics[width=\textwidth]{field-model-blender.png}
    \caption{Blender дээр тоглох талбайн модельд материал олгосны дараа}
    \label{fig:field-blender}
\end{figure}

\newpage

\subsection{Сургах сан}

Зураг таних машин сургалтад зориулсан Ultralytics YOLO санг ашиглан сургаж, судар таних систем боловсруулна.

Ultralytics YOLO зураг таних систем нь тухайн обьект ямар төрлийн обьект болон камерт байх байршлыг нь тогтоох боломжтой байдаг.

\begin{figure}[h]
    \centering
    
\includegraphics[width=0.5\textwidth]{ultralytics.png}
    \caption{Ultralytics YOLO лого}
    \label{fig:ultralytics}
\end{figure}